
\subsubsection*{Correct Specification of the Error Covariance Matrix}

The OLS regression summary produced by \texttt{statsmodels} may contain the
following note:

\begin{quote}
\texttt{[1] Standard Errors assume that the covariance matrix of the errors
is correctly specified.}
\end{quote}

This note concerns the validity of the \textbf{standard errors} and,
consequently, the validity of the \(t\)-statistics, \(p\)-values, and
confidence intervals reported by the regression.

It does not, by itself, mean that \texttt{statsmodels} has detected a problem.
Rather, it reminds us that statistical inference depends on the covariance
structure assumed for the regression errors.


% ============================================================
% OLS MODEL
% ============================================================

\subsubsection*{The Classical OLS Error Covariance Assumption}

Consider the linear regression model

\[
\boxed{
\mathbf{y}
=
X\boldsymbol{\beta}
+
\boldsymbol{\varepsilon},
}
\]

where

\[
\mathbf{y}\in\mathbb{R}^{n},
\qquad
X\in\mathbb{R}^{n\times k},
\qquad
\boldsymbol{\beta}\in\mathbb{R}^{k},
\qquad
\boldsymbol{\varepsilon}\in\mathbb{R}^{n}.
\]

The error vector is

\[
\boldsymbol{\varepsilon}
=
\begin{pmatrix}
\varepsilon_1\\
\varepsilon_2\\
\vdots\\
\varepsilon_n
\end{pmatrix}.
\]

With the default

\begin{verbatim}
Covariance Type: nonrobust
\end{verbatim}

the classical OLS covariance assumption is

\[
\boxed{
\operatorname{Var}
\left(
\boldsymbol{\varepsilon}\mid X
\right)
=
\sigma^2 I_n.
}
\]

Explicitly,

\[
\sigma^2 I_n
=
\begin{pmatrix}
\sigma^2 & 0 & \cdots & 0\\
0 & \sigma^2 & \cdots & 0\\
\vdots & \vdots & \ddots & \vdots\\
0 & 0 & \cdots & \sigma^2
\end{pmatrix}.
\]

This assumption contains two important components.


% ============================================================
% HOMOSKEDASTICITY
% ============================================================

\paragraph{1. Homoskedasticity}

All observations are assumed to have the same conditional error variance:

\[
\boxed{
\operatorname{Var}
\left(
\varepsilon_i\mid X
\right)
=
\sigma^2
\qquad
\forall i.
}
\]

Thus, the magnitude of the error variance does not systematically change
across observations.


% ============================================================
% NO ERROR CORRELATION
% ============================================================

\paragraph{2. No Correlation Between Errors}

For two different observations,

\[
i\neq j,
\]

the errors are assumed to be conditionally uncorrelated:

\[
\boxed{
\operatorname{Cov}
\left(
\varepsilon_i,\varepsilon_j\mid X
\right)
=
0.
}
\]

For time-series data, this assumption implies, in particular, the absence of
serial correlation such as

\[
\operatorname{Cov}
\left(
\varepsilon_t,\varepsilon_{t-1}\mid X
\right)
=
0.
\]

Therefore,

\[
\boxed{
\operatorname{Var}
(\boldsymbol{\varepsilon}\mid X)
=
\sigma^2I_n
}
\]

combines constant error variance and zero covariance between different
errors.


% ============================================================
% OLS ESTIMATOR
% ============================================================

\subsubsection*{OLS Estimator}

The OLS estimator is

\[
\boxed{
\widehat{\boldsymbol{\beta}}
=
(X^\top X)^{-1}X^\top\mathbf{y}.
}
\]

Substituting

\[
\mathbf{y}
=
X\boldsymbol{\beta}
+
\boldsymbol{\varepsilon},
\]

we obtain

\[
\begin{aligned}
\widehat{\boldsymbol{\beta}}
&=
(X^\top X)^{-1}
X^\top
\left(
X\boldsymbol{\beta}
+
\boldsymbol{\varepsilon}
\right)
\\
&=
(X^\top X)^{-1}X^\top X\boldsymbol{\beta}
+
(X^\top X)^{-1}X^\top\boldsymbol{\varepsilon}
\\
&=
\boldsymbol{\beta}
+
(X^\top X)^{-1}X^\top\boldsymbol{\varepsilon}.
\end{aligned}
\]

Hence,

\[
\boxed{
\widehat{\boldsymbol{\beta}}
-
\boldsymbol{\beta}
=
(X^\top X)^{-1}
X^\top
\boldsymbol{\varepsilon}.
}
\]


% ============================================================
% COVARIANCE OF OLS
% ============================================================

\subsubsection*{Covariance Matrix of the OLS Estimator}

Conditional on \(X\),

\[
\operatorname{Var}
\left(
\widehat{\boldsymbol{\beta}}\mid X
\right)
=
\operatorname{Var}
\left[
(X^\top X)^{-1}
X^\top
\boldsymbol{\varepsilon}
\mid X
\right].
\]

Using the general rule

\[
\operatorname{Var}(A\mathbf{x})
=
A\operatorname{Var}(\mathbf{x})A^\top,
\]

we obtain

\[
\boxed{
\operatorname{Var}
\left(
\widehat{\boldsymbol{\beta}}\mid X
\right)
=
(X^\top X)^{-1}
X^\top
\operatorname{Var}
(\boldsymbol{\varepsilon}\mid X)
X
(X^\top X)^{-1}.
}
\]

This is the \textbf{general expression} for the covariance matrix of the OLS
estimator.


% ============================================================
% CLASSICAL CASE
% ============================================================

\subsubsection*{Classical Homoskedastic Case}

Under the classical assumption,

\[
\operatorname{Var}
(\boldsymbol{\varepsilon}\mid X)
=
\sigma^2I_n.
\]

Substituting into the general formula gives

\[
\begin{aligned}
\operatorname{Var}
\left(
\widehat{\boldsymbol{\beta}}\mid X
\right)
&=
(X^\top X)^{-1}
X^\top
(\sigma^2I_n)
X
(X^\top X)^{-1}
\\
&=
\sigma^2
(X^\top X)^{-1}
X^\top X
(X^\top X)^{-1}
\\
&=
\boxed{
\sigma^2
(X^\top X)^{-1}.
}
\end{aligned}
\]

Since the true error variance \(\sigma^2\) is unknown, it must be estimated.

The usual estimator is

\[
\boxed{
\widehat{\sigma}^2
=
\frac{
\sum_{i=1}^{n}
\widehat{\varepsilon}_i^2
}{
n-k
}.
}
\]

The classical estimated covariance matrix is therefore

\[
\boxed{
\widehat{\operatorname{Var}}
\left(
\widehat{\boldsymbol{\beta}}
\right)
=
\widehat{\sigma}^2
(X^\top X)^{-1}.
}
\]


% ============================================================
% STANDARD ERRORS
% ============================================================

\subsubsection*{From the Covariance Matrix to Standard Errors}

The diagonal elements of

\[
\widehat{\operatorname{Var}}
\left(
\widehat{\boldsymbol{\beta}}
\right)
\]

are the estimated variances of the individual coefficient estimators.

For coefficient \(j\),

\[
\boxed{
\widehat{\operatorname{Var}}
(\widehat{\beta}_j)
=
\left[
\widehat{\operatorname{Var}}
(\widehat{\boldsymbol{\beta}})
\right]_{jj}.
}
\]

The corresponding standard error is

\[
\boxed{
SE(\widehat{\beta}_j)
=
\sqrt{
\left[
\widehat{\operatorname{Var}}
(\widehat{\boldsymbol{\beta}})
\right]_{jj}
}.
}
\]

The standard error is then used to construct the \(t\)-statistic:

\[
\boxed{
t_j
=
\frac{
\widehat{\beta}_j-\beta_{j,0}
}{
SE(\widehat{\beta}_j)
}.
}
\]

For the usual null hypothesis

\[
H_0:\beta_j=0,
\]

this becomes

\[
\boxed{
t_j
=
\frac{
\widehat{\beta}_j
}{
SE(\widehat{\beta}_j)
}.
}
\]

The \(t\)-statistic is subsequently used to compute the \(p\)-value and
confidence interval.

Therefore, there is a direct chain:

\[
\boxed{
\text{error covariance assumption}
\longrightarrow
\widehat{\operatorname{Var}}(\widehat{\boldsymbol{\beta}})
\longrightarrow
SE(\widehat{\beta}_j)
\longrightarrow
t_j
\longrightarrow
p\text{-value and confidence interval}.
}
\]


% ============================================================
% MISSPECIFIED COVARIANCE
% ============================================================

\subsubsection*{What Happens if the Error Covariance Is Misspecified?}

Suppose that the classical assumption

\[
\operatorname{Var}
(\boldsymbol{\varepsilon}\mid X)
=
\sigma^2I_n
\]

is incorrect.

More generally, suppose

\[
\boxed{
\operatorname{Var}
(\boldsymbol{\varepsilon}\mid X)
=
\Omega,
}
\]

where

\[
\Omega
\neq
\sigma^2I_n.
\]

Then the correct covariance matrix of the OLS estimator is

\[
\boxed{
\operatorname{Var}
\left(
\widehat{\boldsymbol{\beta}}\mid X
\right)
=
(X^\top X)^{-1}
X^\top
\Omega
X
(X^\top X)^{-1}.
}
\]

Consequently, the classical formula

\[
\widehat{\sigma}^2(X^\top X)^{-1}
\]

is generally no longer an appropriate estimator of the covariance matrix.


% ============================================================
% HETEROSKEDASTICITY
% ============================================================

\paragraph{Example: Heteroskedasticity}

Suppose that

\[
\operatorname{Var}
(\varepsilon_i\mid X)
=
\sigma_i^2,
\]

so that the error variance differs across observations.

Then

\[
\Omega
=
\begin{pmatrix}
\sigma_1^2 & 0 & \cdots & 0\\
0 & \sigma_2^2 & \cdots & 0\\
\vdots & \vdots & \ddots & \vdots\\
0 & 0 & \cdots & \sigma_n^2
\end{pmatrix}.
\]

Unless

\[
\sigma_1^2
=
\sigma_2^2
=
\cdots
=
\sigma_n^2,
\]

we have

\[
\Omega
\neq
\sigma^2I_n.
\]

Therefore, the conventional nonrobust OLS standard errors may be incorrect.


% ============================================================
% SERIAL CORRELATION
% ============================================================

\paragraph{Example: Serial Correlation}

For time-series data, errors may also be correlated over time:

\[
\operatorname{Cov}
(\varepsilon_t,\varepsilon_{t-h}\mid X)
\neq0.
\]

In that case, \(\Omega\) contains non-zero off-diagonal elements:

\[
\Omega
=
\begin{pmatrix}
\sigma_1^2
&
\operatorname{Cov}(\varepsilon_1,\varepsilon_2)
&
\cdots
\\
\operatorname{Cov}(\varepsilon_2,\varepsilon_1)
&
\sigma_2^2
&
\cdots
\\
\vdots
&
\vdots
&
\ddots
\end{pmatrix}.
\]

Again,

\[
\Omega
\neq
\sigma^2I_n,
\]

so the classical covariance estimator may be inappropriate.


% ============================================================
% COEFFICIENTS VS INFERENCE
% ============================================================

\subsubsection*{Coefficient Estimation versus Statistical Inference}

A crucial distinction must be made between the OLS coefficient estimates and
their estimated standard errors.

Under the exogeneity assumption

\[
\boxed{
\mathbb{E}
[
\boldsymbol{\varepsilon}\mid X
]
=
\mathbf{0},
}
\]

we have

\[
\begin{aligned}
\mathbb{E}
[
\widehat{\boldsymbol{\beta}}\mid X
]
&=
\boldsymbol{\beta}
+
(X^\top X)^{-1}
X^\top
\mathbb{E}
[
\boldsymbol{\varepsilon}\mid X
]
\\
&=
\boldsymbol{\beta}.
\end{aligned}
\]

Therefore,

\[
\boxed{
\mathbb{E}
[
\widehat{\boldsymbol{\beta}}\mid X
]
=
\boldsymbol{\beta}.
}
\]

Thus, heteroskedasticity or serial correlation does not, by itself,
necessarily make the OLS coefficient estimator biased when the relevant
exogeneity conditions hold.

However, the conventional covariance estimator may be incorrect.

Therefore,

\[
\boxed{
\begin{array}{c}
\text{incorrect covariance specification}
\\
\Downarrow
\\
\text{incorrect estimated coefficient variances}
\\
\Downarrow
\\
\text{incorrect standard errors}
\\
\Downarrow
\\
\text{incorrect }t\text{-statistics}
\\
\Downarrow
\\
\text{potentially incorrect }p\text{-values and confidence intervals}.
\end{array}
}
\]

For example, suppose

\[
\widehat{\beta}_j=0.10.
\]

With a classical standard error of

\[
SE_{\mathrm{classical}}
(\widehat{\beta}_j)
=
0.04,
\]

we obtain

\[
t
=
\frac{0.10}{0.04}
=
2.5.
\]

This might suggest statistical significance.

However, suppose a robust covariance estimator gives

\[
SE_{\mathrm{robust}}
(\widehat{\beta}_j)
=
0.07.
\]

Then

\[
t_{\mathrm{robust}}
=
\frac{0.10}{0.07}
\approx1.43.
\]

The estimated coefficient is unchanged:

\[
\boxed{
\widehat{\beta}_j=0.10,
}
\]

but the statistical inference changes because the estimated uncertainty
around the coefficient has changed.


% ============================================================
% ROBUST STANDARD ERRORS
% ============================================================

\subsubsection*{Heteroskedasticity-Robust Standard Errors}

If heteroskedasticity is suspected, one can use a
heteroskedasticity-consistent covariance estimator.

For example, in \texttt{statsmodels},

\begin{verbatim}
model = sm.OLS(y, X).fit(cov_type='HC3')
\end{verbatim}

uses HC3 robust standard errors.

The OLS coefficient estimates remain the same:

\[
\boxed{
\widehat{\boldsymbol{\beta}}_{\mathrm{HC3}}
=
\widehat{\boldsymbol{\beta}}_{\mathrm{OLS}},
}
\]

but the estimated covariance matrix changes:

\[
\boxed{
\widehat{\operatorname{Var}}
_{\mathrm{HC3}}
(\widehat{\boldsymbol{\beta}})
\neq
\widehat{\sigma}^2
(X^\top X)^{-1}.
}
\]

Consequently, the standard errors, \(t\)-statistics, \(p\)-values, and
confidence intervals may change.


% ============================================================
% HAC / NEWEY WEST
% ============================================================

\subsubsection*{HAC / Newey--West Standard Errors}

For time-series regressions, the residuals may exhibit both

\[
\boxed{
\text{heteroskedasticity}
}
\]

and

\[
\boxed{
\text{serial correlation}.
}
\]

In this situation, a Heteroskedasticity and Autocorrelation Consistent
(HAC) covariance estimator, such as Newey--West, can be used.

For example,

\begin{verbatim}
model = sm.OLS(y, X).fit(
    cov_type='HAC',
    cov_kwds={'maxlags': 5}
)
\end{verbatim}

allows the estimated covariance matrix to be robust to both
heteroskedasticity and autocorrelation up to the chosen lag structure.

Again, the OLS coefficient estimates are unchanged, but their estimated
covariance matrix changes.

\subsubsection*{Alternative Least-Squares Estimators}

In this section, we focus primarily on the estimation of linear regression
models using \textbf{Ordinary Least Squares (OLS)}. However, other
least-squares estimators can be used when the classical assumptions on the
error covariance matrix are not appropriate.

\paragraph{Ordinary Least Squares (OLS)}

OLS assumes, in its classical covariance specification,

\[
\operatorname{Var}(\boldsymbol{\varepsilon}\mid X)
=
\sigma^2 I.
\]

It estimates the coefficients by minimizing the unweighted sum of squared
residuals:

\[
\boxed{
\widehat{\boldsymbol{\beta}}_{\mathrm{OLS}}
=
\arg\min_{\boldsymbol{\beta}}
\sum_{i=1}^{n}
\left(
y_i-\mathbf{x}_i^\top\boldsymbol{\beta}
\right)^2.
}
\]

\paragraph{Weighted Least Squares (WLS)}

WLS is useful when observations have different error variances
(\textbf{heteroskedasticity}). Observations are assigned different weights:

\[
\boxed{
\widehat{\boldsymbol{\beta}}_{\mathrm{WLS}}
=
\arg\min_{\boldsymbol{\beta}}
\sum_{i=1}^{n}
w_i
\left(
y_i-\mathbf{x}_i^\top\boldsymbol{\beta}
\right)^2.
}
\]

Typically, observations with higher error variance receive lower weights:

\[
w_i\propto\frac{1}{\sigma_i^2}.
\]

\paragraph{Generalized Least Squares (GLS)}

GLS is the more general case in which

\[
\boxed{
\operatorname{Var}(\boldsymbol{\varepsilon}\mid X)
=
\Omega,
}
\]

where \(\Omega\) may contain both unequal variances and correlations between
errors.

GLS estimates

\[
\boxed{
\widehat{\boldsymbol{\beta}}_{\mathrm{GLS}}
=
(X^\top\Omega^{-1}X)^{-1}
X^\top\Omega^{-1}\mathbf{y}.
}
\]

WLS can therefore be viewed as a special case of GLS in which \(\Omega\) is
diagonal:

\[
\boxed{
\Omega
=
\operatorname{diag}
(\sigma_1^2,\ldots,\sigma_n^2).
}
\]

Hence,

\[
\boxed{
\text{OLS}
\subset
\text{WLS}
\subset
\text{GLS},
}
\]

in the sense that OLS corresponds to equal error variances, WLS allows
observation-specific variances, and GLS additionally allows a general
covariance structure between errors.